% !TeX root = ../../Book2.tex
\chapter{Bessel 函数与径向 Fourier 分析}
\label{b2:app:A}

\begin{chapterabstract}
本附录从级数构造第一类 Bessel 函数，以积分表示推导固定实阶的大参数渐近，再由无穷远归一化构造第二类解。球面角积分把 Bessel 函数带入本卷的径向 Fourier 公式，并在几何阶数上给出 Hankel 变换的等距与反演。最后计算 Bochner--Riesz 核的全部常数，证明其绝对可积性的准确阈值，说明核衰减与乘子有界性之间的区别。
\end{chapterabstract}
\begin{learninggoals}
\begin{itemize}
\item\label{b2:item:A2-goal-bessel}核验 Bessel 级数、递推及两类解的归一化，区分原点与无穷远的选择原则。
\item\label{b2:item:A2-goal-radial}在固定 \(2\pi\) 约定下计算球面测度和径向函数的 Fourier 变换。
\item\label{b2:item:A2-goal-hankel}辨认 Hankel 变换的积分测度，并从径向 Plancherel 推出几何阶数的 \(L^2\) 理论。
\item\label{b2:item:A2-goal-kernel}推导 Bochner--Riesz 核，检查振荡渐近、体积因子与可积性阈值。
\end{itemize}
\end{learninggoals}

球与圆具有旋转对称性，但这种对称并不会把 Fourier 积分变成普通的一维 Fourier 变换：角向变量的积分留下一个新的振荡核。它在二维叫作 \(J_0\)，在更高维出现不同阶数，在三维又简化为正弦除以半径。把它们放在同一个 Bessel 方程中，可以同时解释圆盘模态、球面衰减和球形频率截断。

\ProseParagraphBreak

本附录以实变量为主，先独立证明后续需要的级数、积分与渐近公式，所用基础是 Gamma 正实积分、常微分方程、Fourier 反演及 Plancherel。标准记法与更广泛的复参数公式可查 \cite[Chapter 10]{NationalInstituteofStandardsNIST}；本附录明确限定固定实阶渐近和几何阶数的 Hankel 酉性，不把这些证明误当作任意复阶或阶数同时增长时的统一估计。

% !TeX root = ../../Book2.tex
\section{Bessel 方程及其标准解}
\label{b2:sec:A201}

对径向函数作 Fourier 变换时，角向积分不再是一个指数函数，而会产生随半径振荡的特殊函数。同一函数也出现在圆盘振动的分离变量法中：角向模态固定后，径向方程多出一个反平方势项。Bessel 方程正是连接这两种计算的常微分方程。

\ProseParagraphBreak
先看圆盘上的一个实际计算。设 \(\kappa>0\)，在半径为 \(a>0\) 的圆盘内考虑
\(-\Delta u=\kappa^2u\)。写
\(x=\rho\cos\theta\)、\(y=\rho\sin\theta\)，其中 \(0<\rho<a\)。
链式法则给出
\[
 \partial_x=\cos\theta\,\partial_\rho
 -\frac{\sin\theta}{\rho}\partial_\theta,\quad
 \partial_y=\sin\theta\,\partial_\rho
 +\frac{\cos\theta}{\rho}\partial_\theta.
\]
将两个一阶算子各自作用两次并相加，混合导数项抵消，
而对系数求导留下 \(\rho^{-1}\partial_\rho\)，所以
\[
 \Delta=\partial_\rho^2+\rho^{-1}\partial_\rho
 +\rho^{-2}\partial_\theta^2.
\]
固定整数角向模态 \(m\in\mathbb Z\)，取
\(u(\rho,\theta)=R(\rho)e^{\mathrm im\theta}\)；
整数条件保证 \(\theta\) 增加 \(2\pi\) 后函数不变。
由于 \(\partial_\theta^2e^{\mathrm im\theta}=-m^2e^{\mathrm im\theta}\)，原方程化为
\begin{equation}\label{b2:eq:A2-disk-radial}
 \rho^2R''(\rho)+\rho R'(\rho)
 +(\kappa^2\rho^2-m^2)R(\rho)=0.
\end{equation}
再令 \(r=\kappa\rho\)、\(y(r)=R(r/\kappa)\)，便得到
\[
 r^2y''(r)+ry'(r)+(r^2-\nu^2)y(r)=0,
 \quad \nu=|m|.
\]
模态 \(m=0\) 是圆盘上的径向振动，此时 \(\nu=0\)；
其他模态产生 \(\nu^2/r^2\) 项。这个推导同时说明，原点处的适当条件来自圆盘中心的正则性，
不是在奇点上任意给定两个常微分方程初值。

\subsection{级数解与原点行为}
\label{b2:subsec:A20101}

取实阶参数 \(\nu\in\mathbb R\)。对 \(r>0\) 和 \(C^2\) 函数 \(y:(0,\infty)\to\mathbb C\)，\term[bessel-fangcheng]{Bessel 方程}[Bessel equation]为
\[r^2y''+ry'+(r^2-\nu^2)y=0.\]
本附录主要处理实阶参数；先取 \(\nu>-1\)。\chapref{b2:ch:06}已经定义正实参数的 Gamma 积分。分部积分给出 \(\Gamma(s+1)=s\Gamma(s)\)、\(\Gamma(1)=1\)，Gauss 积分给出 \(\Gamma(1/2)=\sqrt\pi\)。这些等式足以处理下面全部实阶级数与几何常数。

\ProseParagraphBreak
在原点附近，暂时略去较低影响的 \(r^2y\) 项，试代 \(y=r^\sigma\)，
得到指标方程 \(\sigma^2-\nu^2=0\)。
为了构造其中的 \(r^\nu\) 分支，进一步寻找
\(y=r^\nu\sum_{k\geq0}c_kr^{2k}\)。
采用偶次步长，是因为方程中的乘法项 \(r^2y\) 恰好把幂次提高二。
算子 \(r^2\partial_r^2+r\partial_r-\nu^2\) 作用于
\(r^{2k+\nu}\) 时给出
\[
 \bigl((2k+\nu)^2-\nu^2\bigr)r^{2k+\nu}
 =4k(k+\nu)r^{2k+\nu}.
\]
于是 \(k=0\) 项自动消失，其他项要求
\[
 c_k=-\frac{c_{k-1}}{4k(k+\nu)}
 =\frac{(-1)^kc_0}
 {4^kk!\,(\nu+1)\cdots(\nu+k)}\quad(k\geq1).
\]
条件 \(\nu>-1\) 保证这些分母不为零。
取 \(c_0=2^{-\nu}/\Gamma(\nu+1)\)，再把连乘积写成 Gamma 比值，
就得到下面的标准归一化。

\Needspace{6\baselineskip}
\begin{definition}[第一类 Bessel 函数]
\label{b2:def:A2-bessel-j}
对 \(\nu>-1\)、\(r>0\)，定义
\[
 J_\nu(r)=\sum_{k=0}^\infty
 \frac{(-1)^k}{k!\,\Gamma(k+\nu+1)}\left(\frac r2\right)^{2k+\nu}.
\]
称它为 \term[diyilei-bessel-hanshu]{第一类 Bessel 函数}[Bessel function of the first kind]。实幂在正半轴取通常正值。
\end{definition}
\symidx{\(J_\nu(r)\)：第一类 Bessel 函数}

\Needspace{6\baselineskip}
\begin{proposition}[级数与微分关系]
\label{b2:prop:A2-bessel-series}
设 \(\nu>-1\)、\(r>0\)，\(J_\nu\) 按\cref{b2:def:A2-bessel-j}的正实幂级数定义。该级数在正半轴每个紧区间上连同任意阶导数一致收敛，并满足 Bessel 方程。原点附近有
\[J_\nu(r)=\frac{(r/2)^\nu}{\Gamma(\nu+1)}(1+O(r^2)).\]
此外，对全部 \(\nu>-1\)、\(r>0\)，有
\[
 \begin{aligned}
 \frac{\mathrm d}{\mathrm dr}(r^{-\nu}J_\nu(r))
 &=-r^{-\nu}J_{\nu+1}(r),\\
 \frac{\mathrm d}{\mathrm dr}(r^{\nu+1}J_{\nu+1}(r))
 &=r^{\nu+1}J_\nu(r).
 \end{aligned}
\]
\end{proposition}
\begin{proof}
相邻项绝对值之比为
\(r^2/[4(k+1)(k+\nu+1)]\)，在任意有界 \(r\) 区间上一致趋零。
更准确地，\(r^{-\nu}J_\nu(r)\) 是关于 \(r^2\) 的收敛半径为无穷的幂级数。
在 \(0<c\leq r\leq C<\infty\) 上，任意固定次微分只增加关于 \(k\) 的多项式因子，
再乘一个由 \(c,C,\nu\) 控制的因子，所以各导数级数也一致收敛。
这既保证逐项微分，也说明结论没有把可能奇异的端点 \(r=0\) 包含进来。
将第 \(k\) 项记为 \(c_kr^{2k+\nu}\)，代入方程后，
\(r^{2k+\nu}\) 的系数恰为
\(4k(k+\nu)c_k+c_{k-1}\)，由 Gamma 递推为零。
首项非零，下一项与首项之比为 \(-r^2/[4(\nu+1)]\)，所以还有较具体的展开
\[
 J_\nu(r)=\frac{(r/2)^\nu}{\Gamma(\nu+1)}
 \left(1-\frac{r^2}{4(\nu+1)}+O(r^4)\right).
\]

第一条导数关系中，\(k=0\) 项求导为零。把其余项的指标写成 \(k=\ell+1\)，得到
\[
 \begin{aligned}
 \frac{\mathrm d}{\mathrm dr}(r^{-\nu}J_\nu(r))
 &=\sum_{k=1}^\infty
 \frac{(-1)^k\,2k\,r^{2k-1}}
 {2^{2k+\nu}k!\,\Gamma(k+\nu+1)}\\
 &=-\sum_{\ell=0}^\infty
 \frac{(-1)^\ell r^{2\ell+1}}
 {2^{2\ell+\nu+1}\ell!\,\Gamma(\ell+\nu+2)}
 =-r^{-\nu}J_{\nu+1}(r).
 \end{aligned}
\]
第二条关系不需变动指标，而是约去 Gamma 递推中的因子：
\[
 \begin{aligned}
 \frac{\mathrm d}{\mathrm dr}(r^{\nu+1}J_{\nu+1}(r))
 &=\sum_{k=0}^\infty
 \frac{(-1)^k(2k+2\nu+2)r^{2k+2\nu+1}}
 {2^{2k+\nu+1}k!\,\Gamma(k+\nu+2)}\\
 &=\sum_{k=0}^\infty
 \frac{(-1)^kr^{2k+2\nu+1}}
 {2^{2k+\nu}k!\,\Gamma(k+\nu+1)}
 =r^{\nu+1}J_\nu(r).
 \end{aligned}
\]
\end{proof}

例如 \(J_0(r)=1-r^2/4+r^4/64-\cdots\)、\(J_1(r)=r/2-r^3/16+\cdots\)。
一般实阶时，光滑的是加权函数 \(r^{-\nu}J_\nu(r)\) 在原点的偶次展开；
不能因此断言 \(J_\nu\) 本身在原点具有任意阶导数。
特别地，\(-1<\nu<0\) 时首项已经发散。
两条导数关系也可改写为
\[
 J_\nu'(r)=\frac{\nu}{r}J_\nu(r)-J_{\nu+1}(r),\quad
 J_{\nu+1}'(r)=J_\nu(r)-\frac{\nu+1}{r}J_{\nu+1}(r).
\]
这些关系把相邻阶数连接起来，但没有改变各公式中 \(\nu>-1\)、\(r>0\) 的范围。

\ProseParagraphBreak
半整数阶给出可以直接检验常数的例子。由
\[
 \Gamma(k+1/2)=\frac{(2k)!}{4^kk!}\sqrt\pi
\]
逐项比较余弦和正弦级数，得到
\[
 \begin{aligned}
 J_{-1/2}(r)
 &=\sqrt{\frac2{\pi r}}
 \sum_{k=0}^\infty\frac{(-1)^kr^{2k}}{(2k)!}
 =\sqrt{\frac2{\pi r}}\cos r,\\
 J_{1/2}(r)
 &=\sqrt{\frac2{\pi r}}
 \sum_{k=0}^\infty\frac{(-1)^kr^{2k+1}}{(2k+1)!}
 =\sqrt{\frac2{\pi r}}\sin r.
 \end{aligned}
\]
它们解释了三维径向 Fourier 公式为何最终只含正弦。一旦阶数不是半整数，保留 Bessel 记号通常比反复展开级数更清楚。

\ProseParagraphBreak

对于圆盘模态 \(\nu=|m|\)，原点有界的这一支与中心正则性相容：
\(J_{|m|}(\kappa\rho)e^{\mathrm im\theta}\) 是
\(\rho^{|m|}e^{\mathrm im\theta}\) 乘以关于 \(\rho^2\) 的收敛级数，
而前者在 \(m\geq0\) 时等于 \((x+\mathrm iy)^m\)，在 \(m<0\) 时等于
\((x-\mathrm iy)^{|m|}\)。
因此这一模态确实光滑地穿过圆心；\(m=0\) 时级数还给出径向导数在中心为零。
方程在 \(r=0\) 有奇点，不能像普通初值问题那样任意指定 \(y(0),y'(0)\)。
在每个离开零点的区间上，二阶常微分方程仍具有通常的两维解空间；
另一支将在大参数归一化后构造，届时才能说明原点有界性为何排除它。

\subsection{积分表示}
\label{b2:subsec:A20102}

对实数 \(a,b>0\)，记 \(B(a,b)=\int_0^1u^{a-1}(1-u)^{b-1}\,\mathrm du\)，称为\term[beta-hanshu]{Beta 函数}[beta function]。两端可积恰好使用了参数正性。
\Needspace{6\baselineskip}
\begin{lemma}[实参数 Beta 积分]\label{b2:lem:A2-beta-integral}
设实数 \(a,b>0\)，
\(B(a,b)=\int_0^1t^{a-1}(1-t)^{b-1}\,\mathrm dt\)，
并对 \(s>0\) 定义 \(\Gamma(s)=\int_0^\infty t^{s-1}e^{-t}\,\mathrm dt\)。
则 \(B(a,b)=\Gamma(a)\Gamma(b)/\Gamma(a+b)\)。
\end{lemma}
\begin{proof}
在第一象限的二重 Gamma 积分中作 \(x=su,y=s(1-u)\)，其中 \(s>0,0<u<1\)，Jacobi 因子为 \(s\)。因此
\[
 \begin{aligned}
 \Gamma(a)\Gamma(b)
 &=\int_0^\infty\int_0^1
 e^{-s}s^{a+b-1}u^{a-1}(1-u)^{b-1}\,\mathrm du\,\mathrm ds\\
 &=\Gamma(a+b)B(a,b).
 \end{aligned}
\]
非负性允许 Tonelli 换序，两个 Gamma 积分有限且 \(\Gamma(a+b)>0\)，相除即得结论。
\end{proof}

\Needspace{6\baselineskip}
\begin{proposition}[Poisson 积分表示]
\label{b2:prop:A2-poisson-bessel}
设实数 \(\nu>-1/2\)、\(r>0\)，\(J_\nu\) 为级数归一化的第一类 Bessel 函数，则
\[
 J_\nu(r)=\frac{(r/2)^\nu}{\sqrt\pi\,\Gamma(\nu+1/2)}
       \int_{-1}^1e^{irt}(1-t^2)^{\nu-1/2}\,\mathrm dt.
\]
\end{proposition}
\begin{proof}
先在有限区间上展开指数函数。所有项的绝对值之和被
\(e^r(1-t^2)^{\nu-1/2}\) 支配，而 \(\nu>-1/2\) 保证这个权在两端可积，
所以可以交换求和与积分。奇数项因对称性消失。
偶数项的积分经 \(u=t^2\) 换元，再用\cref{b2:lem:A2-beta-integral}，得到
\[
 \begin{aligned}
 \int_{-1}^1t^{2k}(1-t^2)^{\nu-1/2}\,\mathrm dt
 &=\int_0^1u^{k-1/2}(1-u)^{\nu-1/2}\,\mathrm du\\
 &=\frac{\Gamma(k+1/2)\Gamma(\nu+1/2)}
 {\Gamma(k+\nu+1)}.
 \end{aligned}
\]
将它与指数级数的系数 \((-1)^kr^{2k}/(2k)!\) 相乘，
再用 \(\Gamma(k+1/2)/(2k)!=\sqrt\pi/(4^kk!)\)，得到
\[
 \int_{-1}^1e^{\mathrm irt}(1-t^2)^{\nu-1/2}\,\mathrm dt
 =\sqrt\pi\,\Gamma(\nu+1/2)
 \sum_{k=0}^\infty
 \frac{(-1)^k(r/2)^{2k}}{k!\,\Gamma(k+\nu+1)}.
\]
乘以陈述中的前因子，恰好恢复 \(J_\nu\) 的级数。
\end{proof}

奇数部分的积分为零，所以也可以把指数替换成 \(\cos(rt)\)，从而直接看出
\(J_\nu(r)\) 为实数。去掉振荡，再对权作一次 Beta 积分，得到
\[
 |J_\nu(r)|\leq\frac{(r/2)^\nu}{\Gamma(\nu+1)}
 \quad(\nu>-1/2,\ r>0).
\]
但这个界在大 \(r\) 时丢失了振荡抵消，不能用来判断 Fourier 核的衰减。
参数边界 \(\nu=-1/2\) 也不能直接代入该积分，因为权变成不可积的
\((1-t^2)^{-1}\)；这一阶仍可使用前面已经证明的余弦公式。
要读出固定阶的精确衰减，需要分析端点 \(t=\pm1\)；
它们对应球面上法向与频率方向平行的两个驻点。

\subsection{大参数渐近与第二类解}
\label{b2:subsec:A20103}

以下渐近均固定阶数 \(\nu\)，令正实变量 \(r\to\infty\)；常数可以依赖阶数，不能据此处理 \(\nu\) 与 \(r\) 同时增大的情形。下面的端点计算将衰减幂和复相位同时确定。它不是对一般振荡积分直接套用绝对值界。

\Needspace{6\baselineskip}
\begin{lemma}[幂型端点的首项]\label{b2:lem:A2-endpoint-asymptotic}
设实数 \(a>-1\)，复函数 \(b\) 在 \([0,\infty)\) 光滑且紧支撑，
其中原点处的光滑性指存在向一个开邻域的光滑延拓。
当正实数 \(r\to+\infty\) 时，
\[
 \int_0^\infty e^{-\mathrm irt}t^ab(t)\,\mathrm dt
 =\Gamma(a+1)e^{-\mathrm i\pi(a+1)/2}b(0)r^{-a-1}
 +O(r^{-a-2}),
\]
余项常数可以依赖固定的 \(a,b\)。
\end{lemma}
\begin{proof}
先计算带指数衰减的模型。设
\(F(r)=\int_0^\infty t^ae^{-(1+\mathrm ir)t}\,\mathrm dt\)。
微分及分部积分分别给出
\[
 \begin{aligned}
 F'(r)&=-\mathrm i\int_0^\infty
 t^{a+1}e^{-(1+\mathrm ir)t}\,\mathrm dt,\\
 (1+\mathrm ir)\int_0^\infty
 t^{a+1}e^{-(1+\mathrm ir)t}\,\mathrm dt&=(a+1)F(r).
 \end{aligned}
\]
故 \(F'(r)=-\mathrm i(a+1)F(r)/(1+\mathrm ir)\)，而
\(F(0)=\Gamma(a+1)\)。解这个实参数微分方程，得到
\[
 \begin{aligned}
 F(r)&=\Gamma(a+1)(1+\mathrm ir)^{-a-1}\\
 &=\Gamma(a+1)e^{-\mathrm i\pi(a+1)/2}r^{-a-1}
 +O(r^{-a-2}),
 \end{aligned}
\]
幂取右半平面的连续辐角。微分时 \(t^{a+1}e^{-t}\) 可积；分部积分时 \(t^{a+1}\) 在原点趋零，所以两步均合法。

写 \(b(t)=b(0)e^{-t}+t c(t)\)。分子在0为零，故 \(c\) 在0光滑；
例如可将差商写成
\(\int_0^1[b'(st)+b(0)e^{-st}]\,\mathrm ds\)，看出各阶导数连续。
在无穷远，\(c\) 及其各阶导数指数衰减。
只需证明余项
\(\int_0^\infty e^{-\mathrm irt}t^{a+1}c(t)\,\mathrm dt=O(r^{-a-2})\)。
远离0的部分是光滑快速衰减函数，可分部积分任意多次。
在0附近先取一个于 \(t\leq2/r\) 支撑的光滑截断，
其直接绝对积分由
\(\int_0^{2/r}t^{a+1}\,\mathrm dt\) 控制，不超过 \(Cr^{-a-2}\)。
其余部分用光滑分割放在
\(t\asymp2^j/r\)、\(j\ge0\) 的环带，直到一个固定小尺度。对每个环带分部积分整数 \(N>a+2\) 次；振幅第 \(N\) 阶导数的积分范数至多
\(C(2^j/r)^{a+2-N}\)，包括截断导数在内。
具体地，若环带尺度为 \(\lambda=2^j/r\)，对 \(t^{a+1}\) 求 \(l\) 次导数产生
\(\lambda^{a+1-l}\)，剩余截断导数至多产生 \(\lambda^{-(N-l)}\)，
而积分区间的长度为 \(O(\lambda)\)，合起来就是所列幂次。
该项因此不超过
\[
Cr^{-N}(2^j/r)^{a+2-N}
=Cr^{-a-2}2^{j(a+2-N)}.
\]
指数严格为负，求几何和给 \(Cr^{-a-2}\)，常数与 \(r\) 无关。各环带截断紧支撑，没有边界项。模型加余项便证明公式；证明还显示常数只用到振幅有限个导数及尾部积分界。
\end{proof}

\Needspace{6\baselineskip}
\begin{theorem}[实轴 Bessel 渐近]
\label{b2:thm:A2-bessel-asymptotic}
设 \(J_\nu\) 为级数归一化的第一类 Bessel 函数。
对固定实数 \(\nu>-1/2\)，令
\(\omega_\nu(r)=r-\pi\nu/2-\pi/4\)。
当正实数 \(r\to+\infty\) 时，
\[
 \begin{aligned}
 J_\nu(r)&=\sqrt{\frac2{\pi r}}\cos\omega_\nu(r)+O(r^{-3/2}),\\
 J_\nu'(r)&=-\sqrt{\frac2{\pi r}}\sin\omega_\nu(r)+O(r^{-3/2}).
 \end{aligned}
\]
对每个固定实数 \(\nu\geq0\)，存在唯一
\(Y_\nu\in C^2((0,\infty);\mathbb R)\)，满足
\(r^2Y_\nu''+rY_\nu'+(r^2-\nu^2)Y_\nu=0\) 以及
\[
 \begin{aligned}
 Y_\nu(r)&=\sqrt{\frac2{\pi r}}\sin\omega_\nu(r)+O(r^{-3/2}),\\
 Y_\nu'(r)&=\sqrt{\frac2{\pi r}}\cos\omega_\nu(r)+O(r^{-3/2}).
 \end{aligned}
\]
所有余项均只断言固定阶数时成立，其常数允许依赖 \(\nu\)。
\end{theorem}
\begin{proof}
写 \(\alpha=\nu-1/2>-1\)，在 Poisson 表示的两个端点附近作对称的光滑分割。
两端之外的振幅光滑且紧支撑于 \((-1,1)\)，反复分部积分给出
\(O(r^{-N})\)，其中 \(N\) 可任意大。
在 \(t=1\) 附近令 \(s=1-t\)，则
\[
 (1-t^2)^\alpha=s^\alpha(2-s)^\alpha,\quad
 e^{\mathrm irt}=e^{\mathrm ir}e^{-\mathrm irs}.
\]
选取等于一的端点截断后，光滑因子在 \(s=0\) 的值为 \(2^\alpha\)。
由\cref{b2:lem:A2-endpoint-asymptotic}，这一端的贡献为
\[
 2^\alpha\Gamma(\alpha+1)r^{-\alpha-1}
 e^{\mathrm i[r-\pi(\alpha+1)/2]}+O(r^{-\alpha-2}).
\]
另一端贡献为其共轭。两端相加得到
\[
 2^{\alpha+1}\Gamma(\alpha+1)r^{-\alpha-1}
 \cos\bigl(r-\pi(\alpha+1)/2\bigr)+O(r^{-\alpha-2}).
\]
选择内部积分的分部积分次数 \(N>\nu+3/2\)，
再乘上 Poisson 前因子
\(r^\nu/[2^\nu\sqrt\pi\,\Gamma(\nu+1/2)]\)。
因为
\[
 \nu-\alpha-1=-\tfrac12,\quad
 2^{\alpha+1-\nu}=\sqrt2,\quad
 \frac{\pi(\alpha+1)}2=\frac{\pi\nu}2+\frac\pi4,
\]
恰好得到 \(J_\nu\) 的余弦主项，端点余项变成 \(O(r^{-3/2})\)。

导数式另行从积分表示求出，不对尚未控制导数的余项作形式微分。
积分对 \(r\) 求导时振幅多出 \(\mathrm it\)，
在右端点取值为 \(\mathrm i\)，在左端点取值为 \(-\mathrm i\)。
因此原来合成 \(2\cos\omega_\nu\) 的两项现在合成
\(-2\sin\omega_\nu\)，同样的端点估计仍有 \(O(r^{-3/2})\) 误差。
前因子求导还产生 \(\nu J_\nu(r)/r=O(r^{-3/2})\)，
故它不改变导数的首项。积分求导合法，因为
\(|t|\cdot(1-t^2)^{\nu-1/2}\) 可积。

为构造第二个解，令 \(w=\sqrt r\,y\)、\(q=\nu^2-1/4\)。方程成为 \(w''+w=q r^{-2}w\)。在充分大的 \([R,\infty)\) 上，考虑
\[
 w(r)=\sqrt{2/\pi}\sin\omega_\nu(r)
 +\int_r^\infty\sin(t-r)\frac q{t^2}w(t)\,\mathrm dt.
\]
自由项选择正弦，是为了确定第二支与 \(J_\nu\) 相差四分之一周期的相位；
积分项补偿方程中尚未包含的 \(q/r^2\) 势。
记积分算子为 \(T\)。在有界连续函数空间 \(C_b([R,\infty))\) 中，
\[
 |T\varphi(r)|\leq
 \int_r^\infty\frac{|q|}{t^2}\cdot|\varphi(t)|\,\mathrm dt
 \leq\frac{|q|}{r}\cdot\|\varphi\|_\infty.
\]
取 \(R>\max\{1,2|q|\}\)，便有 \(\|T\|\leq|q|/R<1/2\)。
收缩原理给出唯一有界连续解 \(w\)，且
\(\|w\|_\infty\leq2\sqrt{2/\pi}\)。
自由项和积分核为实值，所以唯一解也为实值。
对积分项求导时，下端点的正弦为零，因此
\[
 (T\varphi)'(r)
 =-\int_r^\infty\cos(t-r)\frac q{t^2}\varphi(t)\,\mathrm dt.
\]
两项尾积分均为 \(O(r^{-1})\)，故
\[
 \begin{aligned}
 w(r)&=\sqrt{2/\pi}\sin\omega_\nu(r)+O(r^{-1}),\\
 w'(r)&=\sqrt{2/\pi}\cos\omega_\nu(r)+O(r^{-1}).
 \end{aligned}
\]
再求一次导数，连续性和可积尾界给出
\((T\varphi)''=qr^{-2}\varphi-T\varphi\)。
所以 \(w\) 确实满足 \(w''+w=qr^{-2}w\)。
令 \(y=w/\sqrt r\)，并用
\[
 y'(r)=r^{-1/2}w'(r)-\tfrac12r^{-3/2}w(r),
\]
就得到所列 \(y,y'\) 的两项渐近。
取某个 \(r>R\) 上的值与导数作为初值，线性常微分方程的存在唯一性将解延拓到每个
\([\delta,\infty)\)、\(\delta>0\)，因而得到整个正半轴上的解。

若有另一个解满足同样两项渐近，两者之差对应的加权函数 \(v,v'\) 在无穷远趋零。
对等价的一阶系统使用\cref{b2:thm:ode-variation-constants}的变常数公式，从终点 \(L>r\) 向 \(r\) 回溯，得到
\[
 \begin{aligned}
 v(r)={}&v(L)\cos(L-r)-v'(L)\sin(L-r)\\
 &+\int_r^L\sin(t-r)\frac q{t^2}v(t)\,\mathrm dt.
 \end{aligned}
\]
令 \(L\to\infty\)，两项边界值趋零；
由于 \(v\) 在大半轴有界，积分可按 \(t^{-2}\) 支配收敛。
于是得到不含自由项的方程 \(v=Tv\)。
收缩性迫使 \(v\) 在大半轴为零，再由初值唯一性处处为零。
\end{proof}

对实数 \(\nu\geq0\)，将\cref{b2:thm:A2-bessel-asymptotic}中由正弦相位及其导数渐近唯一确定的实 Bessel 解，
称为 \term[dierlei-bessel-hanshu]{第二类 Bessel 函数}[Bessel function of the second kind]，记为 \(Y_\nu\)。
这一归一化确定了正实轴上的第二类解，并与常用约定一致。
定义、积分表示和固定阶渐近的标准记法可对照 NIST《Digital Library of Mathematical Functions》
\cite[\S\S10.2, 10.9, 10.17]{NationalInstituteofStandardsNIST}。
\symidx{\(Y_\nu(r)\)：第二类 Bessel 函数}

\Needspace{6\baselineskip}
\begin{proposition}[Wronski 行列式与原点的第二支]
\label{b2:prop:A2-bessel-wronskian}
令 \(J_\nu\) 取\cref{b2:def:A2-bessel-j}的级数归一化、\(Y_\nu\) 取\cref{b2:thm:A2-bessel-asymptotic}的实轴正弦渐近归一化。对 \(\nu\ge0\)、\(r>0\)，有
\[J_\nu(r)Y_\nu'(r)-J_\nu'(r)Y_\nu(r)=\frac2{\pi r}.\]
因此二者线性无关。原点附近，\(\nu>0\) 时
\[Y_\nu(r)=-\frac{\Gamma(\nu)}\pi\left(\frac2r\right)^\nu(1+o(1)),\]
而 \(Y_0(r)=\frac2\pi\log r+O(1)\)。
\end{proposition}
\begin{proof}
对两解的 Wronski 行列式
\(W=J_\nu Y_\nu'-J_\nu'Y_\nu\) 求导，
交叉项抵消，剩下 \(W'=J_\nu Y_\nu''-J_\nu''Y_\nu\)。
由方程
\(y''=-y'/r-(1-\nu^2/r^2)y\)，可得 \(W'=-W/r\)，
所以 \(rW\) 为常数。用无穷远渐近计算，
\[
 W(r)=\frac2{\pi r}
 \bigl(\cos^2\omega_\nu(r)+\sin^2\omega_\nu(r)\bigr)
 +O(r^{-2}),
\]
故该常数为 \(2/\pi\)，特别地两解线性无关。

由原点级数，存在 \(r_0>0\)，使 \(J_\nu\) 在 \((0,r_0]\) 不为零。
所以在这个区间上
\[
 \left(\frac{Y_\nu}{J_\nu}\right)'=\frac2{\pi rJ_\nu(r)^2}.
\]
令 \(a=2^{-\nu}/\Gamma(\nu+1)>0\)，使用
\(J_\nu(r)=ar^\nu(1+O(r^2))\)，并从 \(r_0\) 积到 \(r\)，得到
\[
 \frac{Y_\nu(r)}{J_\nu(r)}
 =\frac{Y_\nu(r_0)}{J_\nu(r_0)}
 -\frac2{\pi a^2}\int_r^{r_0}s^{-2\nu-1}(1+O(s^2))\,\mathrm ds.
\]
若 \(\nu>0\)，主积分等于
\((r^{-2\nu}-r_0^{-2\nu})/(2\nu)\)。
误差积分由 \(\int_r^{r_0}s^{1-2\nu}\,\mathrm ds\) 控制，因而
\[
 \frac{Y_\nu(r)}{J_\nu(r)}
 =-\frac{r^{-2\nu}}{\pi\nu a^2}
 +\begin{cases}
 O(1),&0<\nu<1,\\
 O(1+|\log r|),&\nu=1,\\
 O(r^{2-2\nu}),&\nu>1.
 \end{cases}
\]
各段余项均为 \(o(r^{-2\nu})\)。
乘回 \(J_\nu\)，再用 \(\Gamma(\nu+1)=\nu\Gamma(\nu)\)，就得到
\[
 -\frac1{\pi\nu a}r^{-\nu}
 =-\frac{\Gamma(\nu)}\pi\left(\frac2r\right)^\nu
\]
作为 \(Y_\nu\) 的首项；中间阶数产生的对数没有改变这一主项。
若 \(\nu=0\)，则 \(a=1\)，被积函数为
\(2/(\pi r)+O(r)\)。
因此 \(Y_0/J_0=(2/\pi)\log r+C+O(r^2)\)。
再乘 \(J_0=1+O(r^2)\)，得到
\(Y_0(r)=(2/\pi)\log r+C+O(r^2|\log r|)\)，特别地有陈述中的有界余项。
\end{proof}

Wronski 常数确定了奇异首项的系数，但仅靠它和原点首项还不能唯一确定第二类解：
\(Y_\nu+cJ_\nu\) 具有同样的 Wronski 行列式和同样的奇异领先项。
无穷远的纯正弦归一化排除了额外的余弦成分，从而固定这个 \(c\)，
也固定了零阶原点展开中的常数 \(C\)。

\ProseParagraphBreak
对 \(\nu\geq0\)，两解组成正半轴上的基本解组。
\(Y_\nu\) 在原点发散，而 \(J_\nu\) 有界，因此任何原点有界的解都是 \(J_\nu\) 的常数倍。
回到\cref{b2:eq:A2-disk-radial}的整数角向模态，这说明中心正则性确实选出前面构造的光滑分支。
无穷远的出入射条件则会选择 \(J_\nu\pm\mathrm iY_\nu\) 的组合，
其首项分别是
\(\sqrt{2/(\pi r)}\,e^{\pm\mathrm i\omega_\nu(r)}\)。
圆盘边值问题与外部散射问题因此不能使用同一条“去掉奇解”的规则。
上述渐近式固定了实阶参数；若阶数与自变量同时增大，还需要控制它们的相对尺度。

% !TeX root = ../../Book2.tex
\section{球面测度与径向 Fourier 变换}
\label{b2:sec:A202}

本卷 Fourier 变换固定为 \(\widehat f(\xi)=\int f(x)e^{-2\pi\mathrm i x\cdot\xi}\,\mathrm dx\)。球面测度采用通常面积测度，既不除以球面积，也不把 \(2\pi\) 吸收到频率中。径向公式中的常数全部由这两项约定决定。把角向变量积掉以后，留下的仍是带有几何权重的积分；本节先计算这个权重，再说明所得一维变换怎样继承全空间的平方范数与反演。

\subsection{球面角积分}
\label{b2:subsec:A20201}

设整数 \(d\ge2\)，\(\sigma\) 为单位球面 \(S^{d-1}\) 的面积测度。旋转不变性说明 \(\widehat\sigma(\xi)\) 只依赖 \(\rho=|\xi|\)，因此可以把频率方向转到第一坐标轴。记 \(\nu=d/2-1\)。这一阶数将由球面的分层面积因子产生。

\Needspace{6\baselineskip}
\begin{theorem}[球面测度的 Fourier 变换]
\label{b2:thm:A2-sphere-fourier}
设整数 \(d\ge2\)，\(\sigma\) 是单位球面 \(S^{d-1}\) 的面积测度，
\(\widehat\sigma(\xi)=\int_{S^{d-1}}e^{-2\pi\mathrm i x\cdot\xi}\,\mathrm d\sigma(x)\)，
\(\rho=|\xi|\)、\(\nu=d/2-1\)。则
\[
\begin{aligned}
 \widehat\sigma(\xi)&=2\pi\rho^{-\nu}J_\nu(2\pi\rho)
 &&(\rho>0),\\
 \widehat\sigma(0)&=|S^{d-1}|=\frac{2\pi^{d/2}}{\Gamma(d/2)}.
\end{aligned}
\]
等价地，对 \(t>0\)，若指数使用参数 \(t\) 而非 \(2\pi\rho\)，则
\[
 \int_{S^{d-1}}e^{-\mathrm i t\omega_1}\,\mathrm d\sigma(\omega)
 =(2\pi)^{d/2}t^{1-d/2}J_{d/2-1}(t).
\]
\end{theorem}
\begin{proof}
先确定球面积。对任意正整数 \(m\)，将
\(\int_{\mathbb R^m}e^{-|x|^2}\,\mathrm dx\) 分别按直角坐标与极坐标计算，得
\[
\begin{aligned}
 \pi^{m/2}
 &=|S^{m-1}|\int_0^\infty e^{-r^2}r^{m-1}\,\mathrm dr\\
 &=\frac{|S^{m-1}|}{2}\Gamma(m/2).
\end{aligned}
\]
第二步用了 \(u=r^2\)。这里各被积函数非负，乘积积分与极坐标积分均可用 Tonelli 处理。因此
\(|S^{m-1}|=2\pi^{m/2}/\Gamma(m/2)\)；\(m=1\) 时，\(S^0=\{-1,1\}\) 上每个点的测度为 \(1\)，总测度为 \(2\)。

以下计算第一坐标为 \(s\) 的球面薄层。先设 \(d\ge3\)。除去两个极点后，球面可参数化为
\[
 \Phi(s,\eta)=\bigl(s,\sqrt{1-s^2}\,\eta\bigr),
 \quad -1<s<1,\quad \eta\in S^{d-2}.
\]
若 \(e\) 是 \(S^{d-2}\) 上的单位切向量，则
\[
\begin{aligned}
 \partial_s\Phi&=\left(1,-\frac{s}{\sqrt{1-s^2}}\eta\right),
 &|\partial_s\Phi|&=(1-s^2)^{-1/2},\\
 D_\eta\Phi(e)&=\bigl(0,\sqrt{1-s^2}\,e\bigr),
 &|D_\eta\Phi(e)|&=\sqrt{1-s^2}.
\end{aligned}
\]
由于 \(\eta\cdot e=0\)，第一种向量与全部角向切向量正交。在角向取正交单位基后，面积的 Jacobi 因子就是这 \(d-1\) 个长度的乘积：
\[
 (1-s^2)^{-1/2}
 \bigl(\sqrt{1-s^2}\bigr)^{d-2}
 =(1-s^2)^{(d-3)/2}.
\]
因此薄层的面积不只是其横截球面的面积，还要乘上沿球面移动时相对于 \(\mathrm ds\) 的伸长因子。两个极点的面积测度为零，故对非负可测 \(F\)，以及满足右侧绝对可积条件的复值 \(F\)，都有
\[
 \int_{S^{d-1}}F(\omega_1)\,\mathrm d\sigma(\omega)
 =|S^{d-2}|\int_{-1}^1F(s)(1-s^2)^{(d-3)/2}\,\mathrm ds.
\]
\(d=2\) 时，直接使用上、下两个半圆
\(\bigl(s,\pm\sqrt{1-s^2}\bigr)\)。每条弧的长度元均为
\((1-s^2)^{-1/2}\,\mathrm ds\)，相加给出同一公式中的
\(|S^0|=2\)。此时两端的奇性仍可积，并没有因分层而产生原球面上不存在的无穷面积。

取 \(F(s)=e^{-\mathrm i ts}\)。此时
\((d-3)/2=\nu-1/2\)，而 \(\nu=d/2-1>-1/2\)，所以
\cref{b2:prop:A2-poisson-bessel}的 Poisson 积分表示适用。将 \(s\) 换成 \(-s\) 后，指数的负号不影响结果，于是
\[
\begin{aligned}
 \int_{S^{d-1}}e^{-\mathrm i t\omega_1}\,\mathrm d\sigma(\omega)
 &=|S^{d-2}|\sqrt\pi\,\Gamma(\nu+1/2)
       \left(\frac2t\right)^\nu J_\nu(t)\\
 &=(2\pi)^{d/2}t^{-\nu}J_\nu(t).
\end{aligned}
\]
令 \(t=2\pi\rho\)，前因子化为 \(2\pi\rho^{-\nu}\)。在 \(\rho=0\) 处，原积分就是总面积；也可用
\cref{b2:prop:A2-bessel-series}的首项
\(J_\nu(t)=(t/2)^\nu/\Gamma(\nu+1)+O(t^{\nu+2})\)
核对这一极限。
\end{proof}

由\cref{b2:thm:A2-bessel-asymptotic}的固定阶大参数渐近，以及零频率附近的连续性，得到
\[
 |\widehat\sigma(\xi)|
 \le C_d(1+|\xi|)^{-(d-1)/2}.
\]
衰减指数是球面维数的一半，恰与驻相公式中每个非退化切向变量贡献半阶衰减一致。两个驻点分别是面向和背向频率方向的极点，它们叠加成一个余弦，因此不能把模长处处替换成严格正的同阶下界。

\ProseParagraphBreak

二维时公式为 \(2\pi J_0(2\pi\rho)\)，三维时为
\(2\sin(2\pi\rho)/\rho\)。三维在 \(\rho=0\) 的极限是 \(4\pi\)，可用来检查面积归一化。半径 \(R>0\) 的球面从单位球面伸缩而来，切空间有 \(d-1\) 个方向，所以其面积测度 \(\sigma_R\) 满足
\[
 \widehat\sigma_R(\xi)=R^{d-1}\widehat\sigma(R\xi).
\]
若使用该球面上的概率测度，还须将右侧除以
\(|S^{d-1}|R^{d-1}\)。这个因子来自测度的归一化，与 Fourier 指数中的 \(2\pi\) 是两件事。

\subsection{径向函数的变换}
\label{b2:subsec:A20202}

对径向函数 \(f(x)=f_0(|x|)\)，极坐标把空间积分分为半径积分与角积分。上一小节已经算出角积分；要把它代入，还须说明两重积分可交换，以及合并为 Bessel 核后采用什么收敛意义。

\Needspace{6\baselineskip}
\begin{theorem}[径向 Fourier 公式]
\label{b2:thm:A2-radial-fourier}
设整数 \(d\ge2\)，\(f_0:(0,\infty)\to\mathbb C\) 可测，
\(f(x)=f_0(|x|)\in L^1(\mathbb R^d)\)，并记
\(\nu=d/2-1\)。则 \(\widehat f(\xi)=F_0(|\xi|)\)，且对 \(\rho>0\)，
\[
 F_0(\rho)=2\pi\rho^{-\nu}\int_0^\infty
       f_0(r)J_\nu(2\pi r\rho)r^{\nu+1}\,\mathrm dr.
\]
该积分绝对收敛。在 \(\rho=0\) 处，
\[
 F_0(0)=|S^{d-1}|\int_0^\infty f_0(r)r^{d-1}\,\mathrm dr,
\]
并且正频率处的公式连续趋于此值。
\end{theorem}
\begin{proof}
先对任意正交变换 \(Q\) 作 \(x=Qy\) 换元。其体积 Jacobi 因子为 \(1\)，且 \(f(Qy)=f(y)\)，所以
\(\widehat f(Q\xi)=\widehat f(\xi)\)。这给出变换的径向性。又由极坐标，
\[
 |S^{d-1}|\int_0^\infty|f_0(r)|r^{d-1}\,\mathrm dr
 =\|f\|_1<\infty.
\]
它正是下面两重积分的绝对积分，因此 Fubini 可用：
\[
 \widehat f(\xi)
 =\int_0^\infty f_0(r)r^{d-1}
    \left(\int_{S^{d-1}}e^{-2\pi\mathrm i r\omega\cdot\xi}
                 \,\mathrm d\sigma(\omega)\right)\mathrm dr.
\]
代入\cref{b2:thm:A2-sphere-fourier}，内层为
\(2\pi(r\rho)^{-\nu}J_\nu(2\pi r\rho)\)。
由于 \(d-1-\nu=\nu+1\)，得到所列公式。

还可以直接给出合并后被积函数的可积控制。角积分的模不超过
\(|S^{d-1}|\)，故对 \(r,\rho>0\)，
\[
 2\pi\rho^{-\nu}|J_\nu(2\pi r\rho)|r^{\nu+1}
 \le |S^{d-1}|r^{d-1}.
\]
乘以 \(|f_0(r)|\) 后右侧可积。这说明公式不是仅靠正负振荡抵消才收敛。最后让 \(\rho\to0\)：内层角积分趋于球面积，上述同一个控制函数允许支配收敛，因而得到零频率处的连续延伸。
\end{proof}

一维径向函数就是偶函数，其公式为
\[
 \widehat f(\rho)=2\int_0^\infty f_0(r)\cos(2\pi r\rho)\,\mathrm dr.
\]
利用 \(J_{-1/2}(t)=\sqrt{2/(\pi t)}\cos t\)，这也可写成径向 Fourier 公式取 \(\nu=-1/2\) 的形式。在一维，球面 \(S^0\) 由两个点组成，角向积分相应化为这两个点上的求和。

\ProseParagraphBreak

对于平方可积函数，径向性应理解为等价类有一个径向代表。若只在某个球面上改动函数，其全空间等价类没有改变；因而不能用每个半径处的点值来定义径向 \(L^2\) 变换。下面先把径向空间和可用的逼近说明清楚。

\Needspace{6\baselineskip}
\begin{lemma}[径向平方可积空间与光滑逼近]
\label{b2:lem:A2-radial-l2}
设 \(d\ge1\) 为整数，\(c_d=|S^{d-1}|\)，其中 \(c_1=2\)。
记 \(X_d=L^2((0,\infty),r^{d-1}\,\mathrm dr;\mathbb C)\)，并令
\(\mathcal R_d\) 为 \(L^2(\mathbb R^d;\mathbb C)\) 中具有
\(h(|x|)\) 形式代表、且 \(h\) 可测的等价类全体。则映射
\[
 V_d:X_d\longrightarrow\mathcal R_d,\quad
 (V_dh)(x)=c_d^{-1/2}h(|x|)
\]
是酉对应，\(\mathcal R_d\) 是闭子空间。由
\(h\in C_c^\infty(0,\infty)\) 产生的径向光滑紧支集函数在
\(\mathcal R_d\) 中稠密。Plancherel 变换保持 \(\mathcal R_d\)，其在此子空间上的限制为酉算子，平方为恒等。
\end{lemma}
\begin{proof}
极坐标给出
\[
 \|V_dh\|_{L^2(\mathbb R^d)}^2
 =\int_0^\infty|h(r)|^2r^{d-1}\,\mathrm dr.
\]
因而 \(V_d\) 不依赖剖面的代表，且为线性等距映射。每个径向平方可积代表的剖面都属于 \(X_d\)，所以其值域恰为
\(\mathcal R_d\)。若 \(V_dh_n\) 在全空间 \(L^2\) 中收敛，上式说明
\((h_n)\) 在完备空间 \(X_d\) 中为 Cauchy 列；设其极限为 \(h\)，则原极限为 \(V_dh\)。这证明值域闭。

对 \(h\in X_d\)，先截去 \(0<r<1/N\) 与 \(r>N\) 两段。
平方积分的绝对连续性使截断误差趋零。在固定区间
\([1/N,N]\) 的稍大紧区间上，权 \(r^{d-1}\) 有正的下界与有限上界。因此把截断函数补零，再用一维光滑磨光作
\(L^2(\mathrm dr)\) 逼近，也得到带权 \(L^2\) 逼近；磨光半径可小到使支集仍离开 \(0\)。逐次选取逼近误差趋零，就得到
\(C_c^\infty(0,\infty)\) 在 \(X_d\) 中的稠密性。相应的
\(h(|x|)\) 在原点邻域恒为零，在其他点则是光滑复合，所以属于
\(C_c^\infty(\mathbb R^d)\)。

这些光滑径向函数的 Fourier 变换由正交换元保持径向性。
由\cref{b2:thm:04-plancherel-extension}的等距连续性，及刚证明的径向子空间闭性，一般径向 \(L^2\) 函数的变换仍在
\(\mathcal R_d\)。\cref{b2:thm:04-unitary}说明全空间变换的平方是反射 \(f(x)\mapsto f(-x)\)，而反射在径向子空间上为恒等。故此限制的平方为恒等，特别它满射，且保持平方范数。
\end{proof}

于是一般 \(L^2\) 径向变换的含义是：取上述径向光滑逼近，先计算每个逼近的变换，再取带权平方范数极限。这个定义独立于逼近列，但没有声称原来的 Bessel 积分对每个频率都收敛。

\ProseParagraphBreak

例如 \(f_0(r)=e^{-\pi a r^2}\)、\(a>0\)，
\cref{b2:thm:03-gaussian}的全空间 Gauss 公式给出
\(F_0(\rho)=a^{-d/2}e^{-\pi\rho^2/a}\)。将它代回径向公式，可以核对相应 Bessel--Gauss 积分中的常数。这里使用的是已知 Fourier 变换，关于一般实阶带幂 Gauss 积分的直接级数计算留作章末训练。

\subsection{几何阶数的 Hankel 变换、等距与反演}
\label{b2:subsec:A20203}

径向公式同时含有输入的幂权与输出的幂权。令
\(g(r)=r^\nu f_0(r)\)，可以把积分中的测度统一为 \(r\,\mathrm dr\)。
对实数 \(\nu>-1\)、可测 \(g:(0,\infty)\to\mathbb C\) 及 \(s>0\)，在
\(\int_0^\infty|g(r)|\cdot|J_\nu(rs)|r\,\mathrm dr<\infty\) 时定义
\term[hankel-bianhuan]{Hankel 变换}[Hankel transform]
\[
 (\mathcal H_\nu g)(s)
 \coloneqq\int_0^\infty g(r)J_\nu(rs)r\,\mathrm dr.
\]
\symidx{\(\mathcal H_\nu\)：以 \(r\,\mathrm dr\) 为测度的 Hankel 变换}
本节在几何阶数 \(\nu=d/2-1\)、整数 \(d\ge1\) 时证明其
\(L^2\) 理论；其他实阶的同类理论需要相应的一维谱论，不由整数维 Fourier 变换自动推出。在这些几何阶数上，径向公式成为
\[
 F_0(\rho)=2\pi\rho^{-\nu}(\mathcal H_\nu g)(2\pi\rho).
\]
输出的 \(2\pi\rho\) 是一个频率伸缩，不能仅从两侧幂权相消就把它省去。

\Needspace{6\baselineskip}
\begin{theorem}[几何阶数的 Hankel 等距与反演]
\label{b2:thm:A2-hankel}
设 \(d\ge1\) 为整数、\(\nu=d/2-1\)。对
\(g\in C_c^\infty(0,\infty;\mathbb C)\)、\(s>0\)，令
\[
 (\mathcal H_\nu g)(s)=\int_0^\infty g(r)J_\nu(rs)r\,\mathrm dr.
\]
则 \(\mathcal H_\nu\) 唯一延拓为
\(L^2((0,\infty),r\,\mathrm dr;\mathbb C)\) 上的酉算子，且
\[
\begin{aligned}
 \|\mathcal H_\nu g\|_{L^2(s\,\mathrm ds)}
 &=\|g\|_{L^2(r\,\mathrm dr)},\\
 \mathcal H_\nu^2g&=g.
\end{aligned}
\]
反演等式对一般 \(g\) 按 \(L^2\) 意义理解。
\end{theorem}
\begin{proof}
先取 \(g\in C_c^\infty(0,\infty)\)。因为其支集离开零点，
\(f_0(r)=r^{-\nu}g(r)\) 产生
\(f(x)=f_0(|x|)\in C_c^\infty(\mathbb R^d)\)。其 Fourier 变换记为
\(F_0(|\xi|)\)。由径向 Plancherel 恒等式消去公共球面积，得
\[
 \int_0^\infty|f_0(r)|^2r^{2\nu+1}\,\mathrm dr
 =\int_0^\infty|F_0(\rho)|^2\rho^{2\nu+1}\,\mathrm d\rho.
\]
左侧等于 \(\int|g(r)|^2r\,\mathrm dr\)。在右侧代入径向公式，逐项保留伸缩常数：
\[
\begin{aligned}
 \int_0^\infty|F_0(\rho)|^2\rho^{2\nu+1}\,\mathrm d\rho
 &=(2\pi)^2\int_0^\infty
       |(\mathcal H_\nu g)(2\pi\rho)|^2\rho\,\mathrm d\rho\\
 &=\int_0^\infty|(\mathcal H_\nu g)(s)|^2s\,\mathrm ds.
\end{aligned}
\]
最后一步用了 \(s=2\pi\rho\)，从而
\(\rho\,\mathrm d\rho=s\,\mathrm ds/(2\pi)^2\)。
由此先得等距性。与\cref{b2:lem:A2-radial-l2}中相同的截断、磨光论证适用于权 \(r\)，所以
\(C_c^\infty(0,\infty)\) 在 \(L^2(r\,\mathrm dr)\) 中稠密。
若 \(g_n\to g\) 于该空间，等距性使
\((\mathcal H_\nu g_n)\) 为 Cauchy 列；以其极限定义变换，得到唯一等距延拓。

还须证明反演，不能只由等距性推出满射。仍从上述紧支集 \(g\) 出发。此时 \(f\) 是 Schwartz 函数，
\cref{b2:thm:05-schwartz-invariance}说明其变换也是 Schwartz 函数，尤其可积。\cref{b2:thm:03-inversion}给出的全空间反演是绝对收敛积分。因为球面的正、负号角积分相同，将它写成极坐标便得到
\[
 f_0(r)=2\pi r^{-\nu}\int_0^\infty
       F_0(\rho)J_\nu(2\pi r\rho)\rho^{\nu+1}\,\mathrm d\rho
 \quad(r>0).
\]
记 \(h=\mathcal H_\nu g\)，则
\(F_0(\rho)=2\pi\rho^{-\nu}h(2\pi\rho)\)。代入上式并用同一个换元，得到
\[
\begin{aligned}
 g(r)
 &=(2\pi)^2\int_0^\infty
      h(2\pi\rho)J_\nu(2\pi r\rho)\rho\,\mathrm d\rho\\
 &=\int_0^\infty h(s)J_\nu(rs)s\,\mathrm ds.
\end{aligned}
\]
这个积分的绝对收敛来自对可积的 \(F_0(|\xi|)\) 使用径向公式时已经证明的控制，不是预设的 Hankel 反演。

最后要把这个积分与刚定义的 \(L^2\) 延拓识别起来。
函数 \(h\) 在正半轴光滑，并已由等距性属于 \(L^2(s\,\mathrm ds)\)。
取 \(0\le\chi_N\le1\)、\(\chi_N\in C_c^\infty(0,\infty)\)，使其在
\([1/N,N]\) 上为 \(1\)，且逐点趋于 \(1\)。令
\(h_N=\chi_N\cdot h\)，则 \(h_N\to h\) 于 \(L^2(s\,\mathrm ds)\)。
一方面，延拓的连续性给出
\(\mathcal H_\nu h_N\to\mathcal H_\nu h\) 于 \(L^2\)；
另一方面，对每个固定 \(r>0\)，上式的绝对可积控制给出
\[
 (\mathcal H_\nu h_N)(r)
 \longrightarrow\int_0^\infty h(s)J_\nu(rs)s\,\mathrm ds=g(r).
\]
从平方范数收敛列取一个几乎处处收敛子列，即知
\(\mathcal H_\nu h=g\) 作为等价类成立。因此
\(\mathcal H_\nu^2g=g\) 先在稠密类成立，再由连续性推广到全部
\(L^2(r\,\mathrm dr)\)。每个 \(g\) 都是
\(\mathcal H_\nu(\mathcal H_\nu g)\)，故算子满射；复线性等距映射保持内积，因而为酉算子。
\end{proof}

这条证明使用的是先前已独立建立的全空间 Fourier 反演与
Plancherel 定理。两次径向积分的迭代没有被直接交换；若直接把两个
Bessel 核相乘并对中间变量积分，通常得不到可用 Fubini 的绝对可积函数。

\ProseParagraphBreak

标准一维 Fourier 变换使用 Lebesgue 测度
\(\mathrm dr\)，上述 Hankel 变换使用 \(r\,\mathrm dr\)。
若希望改用 \(\mathrm dr\)，要把平方范数中的密度同时分配到输入与输出中。下面把两种写法的对应关系列明。

\Needspace{6\baselineskip}
\begin{proposition}[两种测度下的 Hankel 变换]
\label{b2:prop:A2-hankel-lebesgue}
设 \(d\ge1\) 为整数、\(\nu=d/2-1\)，并令
\(\mathcal H_\nu\) 为\cref{b2:thm:A2-hankel}中在
\(L^2((0,\infty),r\,\mathrm dr;\mathbb C)\) 上的酉变换。映射
\[
 U:L^2(r\,\mathrm dr)\longrightarrow L^2(\mathrm dr),
 \quad (Ug)(r)=\sqrt r\,g(r)
\]
为酉算子。共轭变换
\(\mathcal K_\nu\coloneqq U\mathcal H_\nu U^{-1}\)
在 \(L^2((0,\infty),\mathrm dr;\mathbb C)\) 上酉且平方为恒等。
对 \(u\in C_c^\infty(0,\infty;\mathbb C)\)，其积分公式为
\[
 (\mathcal K_\nu u)(s)
 =\int_0^\infty \sqrt{rs}\,J_\nu(rs)u(r)\,\mathrm dr.
\]
\end{proposition}
\begin{proof}
直接计算得
\(\int|Ug(r)|^2\,\mathrm dr=\int|g(r)|^2r\,\mathrm dr\)。
对任意 \(u\in L^2(\mathrm dr)\)，\(r^{-1/2}u(r)\) 属于
\(L^2(r\,\mathrm dr)\) 且映回 \(u\)，所以 \(U\) 满射。
在紧支集测试函数上，
\[
\begin{aligned}
 (U\mathcal H_\nu U^{-1}u)(s)
 &=\sqrt s\int_0^\infty
       r^{-1/2}u(r)J_\nu(rs)r\,\mathrm dr\\
 &=\int_0^\infty \sqrt{rs}\,J_\nu(rs)u(r)\,\mathrm dr.
\end{aligned}
\]
酉性与平方为恒等由共轭关系得到，再以稠密性延拓到全部空间。
\end{proof}

所以把 \(r\,\mathrm dr\) 改成 \(\mathrm dr\) 时，核也从
\(J_\nu(rs)\) 改为 \(\sqrt{rs}J_\nu(rs)\)。不同参考书中的公式若使用不同测度，须先作上述酉对应，再比较常数。

\ProseParagraphBreak

同一个幂权也解释了径向微分算子。对光滑剖面
\(f_0(r)=r^{-\nu}g(r)\)，逐次微分给出
\[
\begin{aligned}
 f_0'&=r^{-\nu}g'-\nu r^{-\nu-1}g,\\
 f_0''&=r^{-\nu}g''-2\nu r^{-\nu-1}g'
            +\nu(\nu+1)r^{-\nu-2}g.
\end{aligned}
\]
由于 \(d-1=2\nu+1\)，径向负 Laplace 表达式因而为
\[
 -\left(f_0''+\frac{d-1}{r}f_0'\right)
 =r^{-\nu}\left(-g''-\frac1r g'+\frac{\nu^2}{r^2}g\right).
\]
括号内是适配测度 \(r\,\mathrm dr\) 的 Bessel 微分表达式。
Bessel 方程给出它的广义特征函数 \(J_\nu(rs)\) 及特征值 \(s^2\)。
要把这句话变成变换公式，还须实际处理分部积分的边界项。

\Needspace{6\baselineskip}
\begin{proposition}[测试函数上的 Bessel 微分表达式]
\label{b2:prop:A2-hankel-differential}
设 \(d\ge1\) 为整数、\(\nu=d/2-1\)。对
\(g\in C_c^\infty(0,\infty;\mathbb C)\)，令
\[
 L_\nu g\coloneqq -g''-\frac1r g'+\frac{\nu^2}{r^2}g.
\]
则对每个 \(s>0\)，都有绝对收敛的恒等式
\[
 \mathcal H_\nu(L_\nu g)(s)=s^2(\mathcal H_\nu g)(s).
\]
此外，
\[
 \int_0^\infty (L_\nu g)(r)\overline{g(r)}r\,\mathrm dr
 =\int_0^\infty
       \left(|g'(r)|^2r+\frac{\nu^2}{r}|g(r)|^2\right)\mathrm dr
 \ge0.
\]
\end{proposition}
\begin{proof}
固定 \(s>0\)，记 \(j(r)=J_\nu(rs)\)，撇号均表示对 \(r\) 微分。
Bessel 方程经链式法则写成
\[
 -(rj')'+\frac{\nu^2}{r}j=s^2rj.
\]
在任意 \(0<a<b<\infty\) 上，两次分部积分给出
\[
\begin{aligned}
 \int_a^b(L_\nu g)j\,r\,\mathrm dr
 &=\int_a^b\left(-(rg')'+\frac{\nu^2}{r}g\right)j\,\mathrm dr\\
 &=\left[-rg'j\right]_a^b
      +\int_a^b\left(rg'j'+\frac{\nu^2}{r}gj\right)\mathrm dr\\
 &=\left[r(gj'-g'j)\right]_a^b
      +\int_a^b g\left(-(rj')'+\frac{\nu^2}{r}j\right)\mathrm dr\\
 &=\left[r(gj'-g'j)\right]_a^b
      +s^2\int_a^b gj\,r\,\mathrm dr.
\end{aligned}
\]
选取 \(a,b\) 使 \(g\) 的支集位于 \((a,b)\)，则 \(g,g'\) 在两端都为零，边界项消失，两个积分也就是全半轴积分。它们的绝对收敛由紧支集离开零点保证。最后把第一步的 \(j\) 换为
\(\overline g\)，作一次分部积分，即得所列非负二次型恒等式。
\end{proof}

离开紧支集测试函数后，边界项需要另行判断。例如，若
\(f_0\) 来自全空间的光滑径向函数，则原点附近
\(f_0(r)=O(1)\)、\(f_0'(r)=O(r)\)。令 \(g=r^\nu f_0\)，由
Bessel 级数
\[
 j(r)=a_s r^\nu(1+O(r^2)),
 \quad a_s=\frac{(s/2)^\nu}{\Gamma(\nu+1)},
\]
以及逐项微分可知，\(r(gj'-g'j)\) 中两个
\(\nu a_s r^{2\nu}f_0\) 项相消，余下为
\[
 r(gj'-g'j)
 =-a_s r^{2\nu+1}f_0'(r)+O(r^{2\nu+2})
 =O(r^d).
\]
故这类正则剖面在原点的边界项趋零，包括 \(d=1\) 的偶函数情形。无穷远仍须另查；例如径向 Schwartz 函数的剖面快速衰减，而
\(J_\nu(sr)\) 及其 \(r\) 导数只有固定阶的幂型大小，所以边界项也趋零，相关积分均可收敛。这些是具体函数类上的判断，不能代替一般定义域条件。

\ProseParagraphBreak

仅有平方可积性甚至不足以保证原点边界项消失。取 \(d=2\)、\(\nu=0\)，令 \(\chi\) 是正半轴上光滑、在原点邻域为 \(1\)、在充分大处为 \(0\) 的函数，置
\(g(r)=\chi(r)\log r\)。则 \(g\in L^2(r\,\mathrm dr)\)，而
\(L_0g\) 在原点邻域为零、在无穷远也为零，故同样属于此空间。但对固定 \(s>0\)，
\(J_0(sr)=1+O(r^2)\)、\(\dfrac{\mathrm d}{\mathrm dr}J_0(sr)=O(r)\)，从而
\[
 r\bigl(g(r)j'(r)-g'(r)j(r)\bigr)\longrightarrow-1
 \quad(r\downarrow0).
\]
因此不能从 \(g,L_\nu g\in L^2(r\,\mathrm dr)\) 就删掉边界项。
若要讨论非紧支撑函数上的算子实现，还须指定零点、无穷远的边界行为及相应定义域。选择自伴实现时，应把这些端点行为写入相应的定义域条件。

% !TeX root = ../../Book2.tex
\section{Bochner--Riesz 核与衰减阈值}
\label{b2:sec:A203}

把频率限制在一个球内，会在物理空间产生长尾振荡；在频率边界加入消失因子，可以改善尾部衰减。Bessel 函数把这种关系写成精确公式。本节只处理核的计算及可积性，精细的 \(L^p\) 乘子范围仍属于第二十四章的问题。

\subsection{球内幂权的变换}
\label{b2:subsec:A20301}

给定 \(\delta>-1\)，定义可积函数
\[m_\delta(\xi)=(1-|\xi|^2)_+^\delta,\]
其中在球外取零，球内取正实幂；球面上的取值不影响积分。\(\delta<0\) 时它在频率边界奇异，所以不能未经检查把它当作有界乘子。其逆 Fourier 变换记为 \(K_\delta\)。

\begin{theorem}[Bochner--Riesz 核]
\label{b2:thm:A2-bochner-riesz}
设整数 \(d\ge1\)、\(\delta>-1\)。对 \(x\ne0\)，有
\[
 K_\delta(x)=\frac{\Gamma(\delta+1)}{\pi^\delta}
       |x|^{-d/2-\delta}J_{d/2+\delta}(2\pi|x|).
\]
它连续延拓到原点，且
\[K_\delta(0)=\frac{\pi^{d/2}\Gamma(\delta+1)}{\Gamma(d/2+\delta+1)}.\]
若频率半径改成 \(R\)，即乘子为 \((1-|\xi|^2/R^2)_+^\delta\)，则核为 \(R^dK_\delta(Rx)\)。
\end{theorem}
\begin{proof}
记 \(\nu=d/2-1\)。先由 Bessel 级数与 Beta 积分证明
\[
 \int_0^1(1-s^2)^\delta s^{\nu+1}J_\nu(ts)\,\mathrm ds
 =2^\delta\Gamma(\delta+1)t^{-\delta-1}J_{\nu+\delta+1}(t).
\]
具体地，把 \(J_\nu\) 的第 \(k\) 项代入，所需积分为
\[\int_0^1s^{2k+2\nu+1}(1-s^2)^\delta\,\mathrm ds
 =\frac{\Gamma(k+\nu+1)\Gamma(\delta+1)}{2\Gamma(k+\nu+\delta+2)}.\]
它与级数分母的一项相消，整理后正是右边的级数。固定 \(t\) 时，绝对值级数在 \(0\le s\le1\) 由一个可积的 \(s^{2\nu+1}(1-s^2)^\delta\) 倍数控制，故逐项积分合法。

对径向乘子应用 \cref{b2:thm:A2-radial-fourier}，令 \(t=2\pi|x|\)，代入上式。因径向函数为偶函数，正、负号 Fourier 指数产生相同角积分。外部系数化为 \(\Gamma(\delta+1)\pi^{-\delta}|x|^{-d/2-\delta}\)，得到核公式；一维使用前节的 \(J_{-1/2}\) 版本即可。Bessel 原点首项给出 \(K_\delta(0)\)，也等于球内幂权的直接积分。最后在逆变换中作 \(\xi=R\eta\) 换元，得到尺度公式。
\end{proof}

例如 \(d=1,\delta=0\) 时，\(K_0(x)=\sin(2\pi x)/(\pi x)\)，原点值为 \(2\)，正是区间 \([-1,1]\) 的长度。\(d=2,\delta=0\) 时，\(K_0(x)=J_1(2\pi|x|)/|x|\)，原点值为 \(\pi\)。这两个低维例子既检查了 \(2\pi\) 的位置，也检查了 Bessel 阶数必须是 \(d/2+\delta\)，而不是球面公式中的 \(d/2-1\)。

\subsection{绝对可积性的准确界限}
\label{b2:subsec:A20302}

\begin{proposition}[核的可积阈值]
\label{b2:prop:A2-kernel-integrability}
对 \(d\ge1\)、\(\delta>-1\)，有
\[K_\delta\in L^1(\mathbb R^d)\quad\Longleftrightarrow\quad
 \delta>\frac{d-1}{2}.\]
更一般地，在无穷远处其领先项为
\[
 \frac{\Gamma(\delta+1)}{\pi^{\delta+1}}
 r^{-(d+1)/2-\delta}
 \cos\left(2\pi r-\frac{\pi(d/2+\delta)}2-\frac\pi4\right),
\]
余项为 \(O(r^{-(d+3)/2-\delta})\)。
\end{proposition}
\begin{proof}
\(d/2+\delta>-1/2\)，故 Bessel 渐近适用，给出所列展开。原点处核连续，无局部可积性障碍；无穷远的绝对积分由 \(\int_1^\infty r^{(d-3)/2-\delta}\,\mathrm dr\) 控制。当 \(\delta>(d-1)/2\) 时，该积分有限。

反之，在每个充分远的单位长度区间内选一个固定长度的子区间，使余弦绝对值至少为 \(3/4\)。余项比主振幅小一个 \(r^{-1}\) 因子，所以在这些子区间上，\(|K_\delta(r)|\ge c r^{-(d+1)/2-\delta}\)。各球壳的绝对积分下界与级数 \(\sum_{n\ge n_0}n^{(d-3)/2-\delta}\) 比较；当 \(\delta\le(d-1)/2\) 时发散。端点相当于调和级数，也不例外。
\end{proof}

频率边界每增加一阶消失，就使物理核多衰减一阶；球面曲率提供的基本衰减则是 \((d-1)/2\) 阶。两者相加后，还要与物理空间的体积增长比较，才能判断绝对可积性。只观察一条径向曲线上的函数值而忘记 \(r^{d-1}\) 体积因子，会得到错误阈值。

\subsection{可积核与乘子问题的区别}
\label{b2:subsec:A20303}

若 \(\delta>(d-1)/2\)，Young 不等式给出 \(f\mapsto K_\delta*f\) 在全部 \(L^p\)、\(1\le p\le\infty\) 上有界，范数不超过 \(\lVert K_\delta\rVert_1\)。缩放后 \(\lVert R^dK_\delta(R\cdot)\rVert_1\) 不变。由于核的总积分为 \(m_\delta(0)=1\)，且绝对积分尾部在缩放下消失，它们构成近似恒等族；于是当 \(R\to\infty\) 时，对 \(1\le p<\infty\)，
\[R^dK_\delta(R\cdot)*f\longrightarrow f\quad\text{于 }L^p.
\]
总积分等式可由 \(m_\delta,K_\delta\in L^1\) 的 Fourier 反演在零点得到，\(m_\delta\) 在零点连续。近似恒等的收敛证明则复用本卷前面已经建立的平移连续性和尾部拆分。

\ProseParagraphBreak

若 \(\delta\) 低于这个阈值，失效的是绝对值积分法，不是所有 \(L^p\) 有界性。例如只要 \(\delta\ge0\)，乘子 \(m_\delta\) 有界，Plancherel 立即给出 \(L^2\) 算子范数至多为 \(1\)，即使对应核不属于 \(L^1\)。核中的振荡抵消仍可能产生额外有界性；第二十四章正是研究这些不能由 Young 不等式看出的范围。

\ProseParagraphBreak

因此有三个不能混用的判断：核是否逐点衰减，核是否绝对可积，以及振荡卷积算子在某个 \(L^p\) 上是否有界。Bessel 公式解决前两个问题并提供第三个问题的必要线索，却不能单独替代限制理论、插值和几何反例。


\begin{chaptersummary}
Bessel 级数确定原点附近的正则支，积分表示把其大参数行为化为两个端点的振荡叠加；第二类解由无穷远的正弦相位归一化，并与第一类解组成基本解组。球面角积分固定了径向 Fourier 公式中的阶数和常数，Hankel 变换则是在适当幂权与测度下的一维表达。

\ProseParagraphBreak

Bochner--Riesz 核的衰减同时记录球面曲率和频率边界的消失阶。绝对可积时可直接用 Young 不等式与近似恒等理论；不可积时仍可能依靠振荡得到 \(L^p\) 有界性。判断具体结论时，应把归一化、积分测度、参数范围和收敛意义一并保留。
\end{chaptersummary}
\BookExercises

下列题目均为本书自拟，围绕具体积分、零点、径向维数变化及演化问题组织。

% 自拟：渐近给出无穷零点，ODE排除重零点。
\begin{exercise}\label{b2:ex:A2-zeros}
设 \(\nu\ge0\)。证明 \(J_\nu\) 有无穷多个正零点，每个正零点均为单零点，而且 \(J_\nu\) 与 \(J_{\nu+1}\) 不共享正零点。
\end{exercise}
\begin{solution}
在余弦相位等于整数倍 \(\pi\) 的充分大点上，渐近主项交替正负，余项比主幅小一个 \(r^{-1}\) 因子，因此每对相邻点间至少有一个零点。若正点同时满足 \(J_\nu=J_\nu'=0\)，普通初值唯一性使该解在正半轴恒为零，与原点首项矛盾，故零点简单。导数公式给出 \(J_\nu'=\nu J_\nu/r-J_{\nu+1}\)；若两个相邻阶数共享正零点，也会迫使导数为零，矛盾。
\end{solution}

% 自拟：不同Bessel零点的正交和范数，包含参数微分。
\begin{exercise}\label{b2:ex:A2-bessel-orthogonality}
设 \(\nu\ge0\)，\(j\ne k\) 是 \(J_\nu\) 的两个正零点。证明
\[\int_0^1rJ_\nu(jr)J_\nu(kr)\,\mathrm dr=0,\qquad
 \int_0^1rJ_\nu(jr)^2\,\mathrm dr=\tfrac12J_{\nu+1}(j)^2.\]
\end{exercise}
\begin{solution}
令 \(L_\nu v=-(rv')'+\nu^2v/r\)。函数 \(v_\lambda(r)=J_\nu(\lambda r)\) 满足 \(L_\nu v_\lambda=\lambda^2rv_\lambda\)。交叉相乘相减并积分，边界项在 \(r=1\) 由两个零点消失，在 \(r=0\) 由级数首项及下一项消失，得 \((j^2-k^2)\int rv_jv_k=0\)。再对参数 \(\lambda\) 求导，记 \(w=\partial_\lambda v_\lambda\)，则 \(L_\nu w=\lambda^2rw+2\lambda rv_\lambda\)。与原方程作同样的 Green 相减，在 \(\lambda=j\) 时得到 \(2j\int rv_j^2=jJ_\nu'(j)^2\)。前题递推式给出 \(J_\nu'(j)=-J_{\nu+1}(j)\)，即范数公式。参数微分和原点边界项均由局部一致收敛级数保证。
\end{solution}

% 自拟：球面二阶矩与变换的低频展开。
\begin{exercise}\label{b2:ex:A2-sphere-moment}
设 \(\sigma\) 为单位球面面积测度。求 \(\int\omega_i\omega_j\,\mathrm d\sigma\)，并由此证明
\[\widehat\sigma(\xi)=|S^{d-1}|\left(1-\frac{2\pi^2}{d}|\xi|^2+O(|\xi|^4)\right).\]
\end{exercise}
\begin{solution}
反射对称性使不同指标的积分为零，旋转对称性使对角项相等；它们的和是 \(\int|\omega|^2\,\mathrm d\sigma=|S^{d-1}|\)，故每项为球面积除以 \(d\)。将指数展开到三阶，奇次项积分为零，二次项为 \(-2\pi^2\int(\omega\cdot\xi)^2=-2\pi^2|S^{d-1}||\xi|^2/d\)。球面紧致，四阶余项可统一积分。也可用 \(J_{d/2-1}\) 的前两项核对常数。
\end{solution}

% 自拟：Hankel变换下的带幂高斯。
\begin{exercise}\label{b2:ex:A2-gaussian}
设 \(a>0\)、\(\nu>-1\)、\(s>0\)。直接证明
\[\int_0^\infty e^{-ar^2}r^{\nu+1}J_\nu(sr)\,\mathrm dr
 =\frac{s^\nu}{(2a)^{\nu+1}}e^{-s^2/(4a)}.\]
\end{exercise}
\begin{solution}
把级数代入。第 \(k\) 项所需的高斯矩是 \(\tfrac12a^{-k-\nu-1}\Gamma(k+\nu+1)\)，与 Bessel 分母相消。因此积分级数等于 \(s^\nu(2a)^{-\nu-1}\sum_{k\ge0}(-s^2/(4a))^k/k!\)。绝对值积分级数也收敛，因为它对应正指数级数 \(e^{s^2/(4a)}\)，故 Fubini 合法。求和即得结论。
\end{solution}

% 自拟：同一径向剖面在相差二维中的变换。
\begin{exercise}\label{b2:ex:A2-dimension-raising}
设光滑紧支集剖面 \(f_0\) 在各维产生径向光滑函数，记其 \(d\) 维变换为 \(F_d(\rho)\)。对 \(d\ge1\)、\(\rho>0\)，证明
\[F_{d+2}(\rho)=-\frac1{2\pi\rho}F_d'(\rho).\]
\end{exercise}
\begin{solution}
令 \(\nu=d/2-1\)。对径向公式求导，并使用 \(\frac{\mathrm d}{\mathrm d\rho}(\rho^{-\nu}J_\nu(2\pi r\rho))=-2\pi r\rho^{-\nu}J_{\nu+1}(2\pi r\rho)\)，得到 \(F_d'=-4\pi^2\rho^{-\nu}\int f_0(r)J_{\nu+1}(2\pi r\rho)r^{\nu+2}\,\mathrm dr\)。与维数 \(d+2\) 的公式比较即得。紧支撑保证求导与积分交换合法；\(d=1\) 可用半整数表达，结论仍成立。
\end{solution}

% 自拟：球形频率环带的显式核和缩放。
\begin{exercise}\label{b2:ex:A2-shell-kernel}
设 \(0<a<b\)，求乘子 \(\mathbf1_{\{a<|\xi|<b\}}\) 的逆 Fourier 变换，并求其原点值。说明该算子在 \(L^2\) 上为何是正交投影，尽管核会改变符号。
\end{exercise}
\begin{solution}
记 \(K_0\) 为单位频率球核，则所求为 \(b^dK_0(bx)-a^dK_0(ax)\)，即两项已知 Bessel 公式之差。原点值为 \(|B_1|(b^d-a^d)\)。乘子实值且平方等于自身，因此在 Plancherel 识别下对应自伴幂等算子，正是正交投影。核变号与投影性质无关；正交投影的正性指 \(\langle Tf,f\rangle\ge0\)，不是指核逐点非负。
\end{solution}

% 自拟：小球面平均与Laplace的系数。
\begin{exercise}\label{b2:ex:A2-small-sphere-average}
设 \(f\in C^4\) 于 \(x\) 的邻域。证明归一化球面平均满足
\[\frac1{|S^{d-1}|}\int_{S^{d-1}}f(x+R\omega)\,\mathrm d\sigma
 =f(x)+\frac{R^2}{2d}\Delta f(x)+O(R^4).\]
说明在 Fourier 侧相同系数怎样出现。
\end{exercise}
\begin{solution}
对 \(f(x+R\omega)\) 作三阶 Taylor 展开，余项由邻域内四阶导数上界统一控制为 \(O(R^4)\)。奇次项因反射对称消失，二阶项用前题矩公式得到 \(R^2\Delta f/(2d)\)。Fourier 侧的归一化乘子是 \(\widehat\sigma(R\xi)/|S^{d-1}|=1-2\pi^2R^2|\xi|^2/d+O(R^4|\xi|^4)\)，而 \(\Delta\) 的符号为 \(-4\pi^2|\xi|^2\)，两者一致。
\end{solution}

% 自拟：Hankel谱表示中的径向热演化。
\begin{exercise}\label{b2:ex:A2-radial-heat}
设 \(d\ge2\)，\(f\in L^2(\mathbb R^d)\) 径向，\(\nu=d/2-1\)，\(g=r^\nu f_0\)。用 \(\mathcal H_\nu\) 表示初值为 \(f\) 的热演化，并证明其 \(L^2\) 范数不增且当 \(t\downarrow0\) 时回到初值。
\end{exercise}
\begin{solution}
热演化的 Fourier 乘子为 \(e^{-4\pi^2t\rho^2}\)。令 \(s=2\pi\rho\)，径向变换公式给出
\[u_0(t,r)=r^{-\nu}\mathcal H_\nu\bigl(e^{-ts^2}\mathcal H_\nu g(s)\bigr)(r),\]
等式按相应 \(L^2\) 意义理解。Hankel 酉性及 \(|e^{-ts^2}|\le1\) 给出带权径向范数不增，乘以公共球面积后就是全空间 \(L^2\) 收缩。由支配收敛，\((e^{-ts^2}-1)\mathcal H_\nu g\to0\) 于 \(L^2(s\,\mathrm ds)\)，再经逆变换得到初值收敛。
\end{solution}
